\section{Introduction and research scope}

The proposed system targets printed circuit boards approximately $70 \times 70$~mm.
An initial imaging field of view of about $80 \times 80$~mm leaves positioning margin;
this is a design hypothesis rather than a confirmed optical specification.

The system will test controlled 365~nm UVA excitation with visible-band imaging
of fluorescence, subject to verifying the spectral response of the actual conformal
coating. It will assess missing coverage, excessive coverage, keep-out violations
and other defect classes only when those categories can be labelled reproducibly.

\subsection{Research questions}
\begin{enumerate}
\item Under what optical conditions does fluorescence-based image analysis
provide repeatable coating-coverage evidence on the target board?
\item How do classical image processing, learned segmentation, and a
reference-aware hybrid differ in defect-level sensitivity and false acceptance?
\item What image resolution, inference model and hardware are sufficient
to meet the measured production inspection time, not just neural-network latency?
\item How robust are the methods to new PCBs, production lots, lighting drift,
geometric displacement, and coating variation?
\end{enumerate}

\subsection{Candidate contribution}
A modular architecture incorporating flat-field correction, PCB registration,
coating segmentation and required/forbidden-region comparisons will be tested
against independent baselines. Novelty and improvement are \emph{hypotheses};
both require literature verification and quantitative evidence.

\subsection{Data governance}
Only explicitly approved test imagery may be released. Industrial artwork,
traceability data and coating specifications may be confidential. Publications
must distinguish accessible publication/test-board data from restricted factory data.
